% =============================================================================
% MIBEL — Energy Analytics Platform
% Relatório Técnico Final
% MEI / TEAD 2.0 | ESTG-IPP | 2025-2026
% =============================================================================
\documentclass[11pt,a4paper]{article}

\usepackage[utf8]{inputenc}
\usepackage[T1]{fontenc}
\usepackage[portuguese]{babel}
\usepackage[margin=2.2cm]{geometry}
\usepackage{graphicx}
\usepackage{booktabs}
\usepackage{longtable}
\usepackage{array}
\usepackage{xcolor}
\usepackage{listings}
\usepackage{hyperref}
\usepackage{enumitem}
\usepackage{titlesec}
\usepackage{fancyhdr}
\usepackage{microtype}
\usepackage{parskip}
\usepackage{tabularx}
\usepackage{multirow}

% Cores
\definecolor{mibel-blue}{RGB}{15,31,63}
\definecolor{mibel-accent}{RGB}{24,56,119}
\definecolor{code-bg}{RGB}{248,248,248}
\definecolor{pass-green}{RGB}{22,163,74}
\definecolor{fail-red}{RGB}{185,28,28}
\definecolor{warn-orange}{RGB}{234,88,12}

% Configuração de títulos
\titleformat{\section}{\large\bfseries\color{mibel-blue}}{}{0em}{}[\titlerule]
\titleformat{\subsection}{\normalsize\bfseries\color{mibel-accent}}{}{0em}{}
\titleformat{\subsubsection}{\normalsize\itshape}{}{0em}{}

% Cabeçalho / Rodapé
\pagestyle{fancy}
\fancyhf{}
\fancyhead[L]{\small\textcolor{mibel-blue}{\textbf{MIBEL — Energy Analytics Platform}}}
\fancyhead[R]{\small\textcolor{gray}{MEI / TEAD 2.0 | ESTG-IPP}}
\fancyfoot[C]{\small\thepage}
\renewcommand{\headrulewidth}{0.4pt}

% Listings (código inline)
\lstset{
  backgroundcolor=\color{code-bg},
  basicstyle=\ttfamily\footnotesize,
  breaklines=true,
  frame=single,
  rulecolor=\color{gray!30},
  xleftmargin=6pt,
  xrightmargin=6pt,
}

% Hyperref
\hypersetup{
  colorlinks=true,
  linkcolor=mibel-accent,
  urlcolor=mibel-accent,
  pdfauthor={Grupo MIBEL — MEI TEAD 2.0},
  pdftitle={MIBEL Energy Analytics Platform — Relatório Técnico},
}

% Sem indentação, espaçamento entre parágrafos
\setlength{\parindent}{0pt}
\setlength{\parskip}{5pt}

% =============================================================================
\begin{document}
% =============================================================================

% Capa
\begin{titlepage}
  \centering
  \vspace*{2cm}
  {\color{mibel-blue}\rule{\linewidth}{1.5pt}}\\[0.4cm]
  {\Huge\bfseries\color{mibel-blue} MIBEL\\[0.3cm]
   \Large Energy Analytics Platform}\\[0.3cm]
  {\color{mibel-blue}\rule{\linewidth}{1.5pt}}\\[1.5cm]
  {\large\textbf{Relatório Técnico Final}}\\[0.5cm]
  {\large Mestrado em Engenharia Informática}\\
  {\large Tecnologias e Engenharia de Dados (TEAD 2.0)}\\[0.5cm]
  {\large ESTG --- Instituto Politécnico do Porto}\\[1.5cm]
  {\large Número de Aluno: \textbf{8220800}}\\[0.3cm]
  {\large Email: \href{mailto:8220800@estg.ipp.pt}{8220800@estg.ipp.pt}}\\[2cm]
  {\large Ano Letivo 2025--2026}\\[1cm]
  \vfill
  {\small Versão 1.0 --- Maio 2026}
\end{titlepage}

\tableofcontents
\newpage

% =============================================================================
\section{Contexto e Cenário}
% =============================================================================

\subsection{Enquadramento do Mercado Energético Ibérico}

O MIBEL (\textit{Mercado Ibérico de Electricidade}) é o mercado grossista de eletricidade da Península Ibérica. Em Portugal, a gestão da rede de transporte é da responsabilidade da REN (Rede Energética Nacional), enquanto os preços \textit{day-ahead} são determinados diariamente na bolsa de energia OMIE. A volatilidade dos preços, a crescente penetração de produção renovável e a dependência de condições meteorológicas criam um ambiente de elevada complexidade analítica.

\subsection{Cenário Criado pelo Grupo}

O grupo criou um cenário hipotético em que uma equipa de \textbf{Data Engineering e ML} de um comercializador de energia português tem como missão:

\begin{enumerate}[leftmargin=*]
  \item \textbf{Monitorizar} em tempo quase-real o equilíbrio entre produção e consumo nacionais;
  \item \textbf{Estimar} o custo horário da energia consumida com base no preço \textit{day-ahead} OMIE;
  \item \textbf{Prever} a produção elétrica diária e o preço \textit{spot} a partir de variáveis meteorológicas, permitindo ao departamento comercial antecipar decisões de compra/venda de energia.
\end{enumerate}

Para responder a este cenário, foi construída a plataforma \textbf{MIBEL Energy Analytics Platform}: um lakehouse local completo, processando dados históricos desde 2023, com três Data Products independentes, quatro modelos de ML em produção e dashboards operacionais.

\subsection{Fontes de Dados}

\begin{table}[h!]
\centering
\begin{tabularx}{\linewidth}{llX}
\toprule
\textbf{Fonte} & \textbf{Data Product} & \textbf{Descrição} \\
\midrule
REN/ERSE CSV & DP-01 & Produção (DGM + PRE) e consumo nacional a 15 min \\
OMIE CSV & DP-02 & Preços \textit{day-ahead} hora a hora (€/MWh) \\
Open-Meteo API & DP-03 & Dados meteorológicos horários para Portugal continental \\
\bottomrule
\end{tabularx}
\end{table}

% =============================================================================
\section{A) Especificação de Data Products (20\%)}
% =============================================================================

\subsection{Os Três Data Products}

Foram especificados três produtos de dados com contratos formais, SLAs/SLOs e estratégia de versionamento:

\begin{table}[h!]
\centering
\begin{tabularx}{\linewidth}{lllX}
\toprule
\textbf{\#} & \textbf{Nome} & \textbf{Grão} & \textbf{Gold Table} \\
\midrule
DP-01 & \texttt{producao\_consumo} & Horário & \texttt{dp\_energia\_balance\_hourly} \\
DP-02 & \texttt{consumo\_preco} & Horário & \texttt{dp\_energy\_market\_hourly} \\
DP-03 & \texttt{meteo\_producao} & Diário & \texttt{dp\_meteo\_producao\_daily\_features} \\
\bottomrule
\end{tabularx}
\end{table}

\subsection{Contratos e SLAs}

Cada DP define: perguntas analíticas, métricas de negócio, grão e chave primária, schema tipado (v1), regras de qualidade e SLAs de \textit{freshness}/completude. Exemplo para DP-01:

\begin{itemize}[leftmargin=*,noitemsep]
  \item \textit{Freshness}: dados disponíveis até T+30~min após fecho da hora
  \item \textit{Completude}: $\geq 99{,}0\%$ das horas no intervalo esperado
  \item \textit{Unicidade}: chave \texttt{timestamp\_utc} sem duplicados
  \item \textit{Qualidade}: \texttt{flag\_defice} e \texttt{flag\_excedente} mutuamente exclusivos
\end{itemize}

\subsection{Decisão de Impacto: \texttt{feature\_set\_version} no DP-03}

O DP-03 inclui a coluna \texttt{feature\_set\_version} (obrigatória) em todas as linhas da tabela Gold. Esta decisão garante que cada \textit{run} de treino ML referencia explicitamente a versão do conjunto de \textit{features} utilizado, permitindo:
\begin{itemize}[leftmargin=*,noitemsep]
  \item Reprodutibilidade: re-treino com dados históricos da mesma versão produz o mesmo modelo
  \item Rastreabilidade: o registo MLflow aponta para \texttt{feature\_set\_version} e \texttt{dataset=iceberg.gold.<tabela>}
  \item Evolução segura: nova transformação de \textit{feature} implica novo major sem quebrar consumidores anteriores
\end{itemize}

% =============================================================================
\section{B) Implementação do Lakehouse (20\%)}
% =============================================================================

\subsection{Stack Tecnológico}

O lakehouse corre integralmente em Docker Compose com os seguintes serviços:

\begin{table}[h!]
\centering
\begin{tabularx}{\linewidth}{llX}
\toprule
\textbf{Serviço} & \textbf{Tecnologia} & \textbf{Função} \\
\midrule
Object Storage & MinIO & S3-compatível; armazena Parquet Bronze e artefactos ML \\
Metastore & Hive 4.0 + MariaDB & Registo de schemas para catálogos Hive e Iceberg \\
Query Engine & Trino 440 & Federação multi-catálogo; execução de SQL sobre Iceberg \\
ML Tracking & MLflow + Postgres & Registo de experimentos, métricas, modelos e artefactos \\
Dashboards & Grafana 10.4 & Visualização direta sobre Trino (sem ETL adicional) \\
Streaming & Redpanda & Broker Kafka-compatível para ingestão em \textit{stream} \\
\bottomrule
\end{tabularx}
\end{table}

\subsection{Arquitetura Medallion}

\begin{description}[leftmargin=*]
  \item[\textbf{Bronze}] Tabelas Hive \textit{external} sobre Parquet em \texttt{s3://warehouse/bronze/}. Sem transformação --- preserva o dado original para auditoria e reprocessamento.
  \item[\textbf{Silver}] Tabelas Iceberg \textit{managed} com deduplicação, validação de \textit{range} físico (limites Portugal continental) e classificação \texttt{\_quality\_flag}: \texttt{ok} / \texttt{null\_values} / \texttt{out\_of\_range}.
  \item[\textbf{Gold}] Tabelas Iceberg \textit{managed} com joins \textit{cross-DP}, lags temporais, \textit{rolling windows} e flags de negócio. Fonte de verdade para APIs, dashboards e ML.
\end{description}

\subsection{Decisão de Impacto: Hive para Bronze, Iceberg para Silver/Gold}

A separação deliberada entre Hive (Bronze) e Iceberg (Silver/Gold) tem impacto direto na qualidade e performance:

\begin{itemize}[leftmargin=*,noitemsep]
  \item \textbf{Iceberg} oferece ACID, \textit{schema evolution} sem rewrite de ficheiros e \textit{time-travel} --- essencial para auditar reprocessamentos e rollback de transformações incorretas
  \item \textbf{Hive external} para Bronze elimina o custo de gestão ACID em dados raw, que são escritos uma única vez e nunca modificados
  \item \textbf{Trino} separa \textit{compute} de \textit{storage}, permitindo escalar o motor de query independentemente do lakehouse
\end{itemize}

% =============================================================================
\section{C) Pipelines com Orquestração (20\%)}
% =============================================================================

\subsection{Runners Python Idempotentes}

Cada DP tem um orquestrador Python autónomo que executa as seguintes fases em sequência:

\begin{enumerate}[leftmargin=*,noitemsep]
  \item Verificação Docker + \texttt{compose up} (opcional com \texttt{--skip-docker})
  \item \textit{Health check} Trino (\texttt{SELECT 1})
  \item Criação de \textit{venv} local + instalação de dependências Python
  \item DDL: \texttt{CREATE TABLE IF NOT EXISTS} (Bronze / Silver / Gold)
  \item Fetch / upload de dados para MinIO (CSV ou Open-Meteo API)
  \item Bronze INSERT (stage Hive $\rightarrow$ Iceberg managed)
  \item Silver INSERT (deduplicação + validação + \texttt{\_quality\_flag})
  \item Gold INSERT (joins, lags, \textit{rolling windows})
  \item Quality Gate (bloqueante em FAIL --- ver Secção~\ref{sec:quality})
\end{enumerate}

\subsection{DP-02 Streaming --- Pipeline de Ingestão Contínua}

O DP-02 tem duas variantes de pipeline que coexistem e partilham a mesma camada Gold:

\begin{description}[leftmargin=*]
  \item[\textbf{Static (OMIE CSV)}] Dados históricos 2022--2024 carregados via Flyte remoto em K3s; adequado para análise retrospetiva e treino ML.
  \item[\textbf{Streaming (ENTSO-E API)}] Ingestão contínua via \textit{ENTSO-E Transparency Platform} (ou Energy-Charts como fallback sem token); atualiza a janela dos últimos 12 meses por defeito.
\end{description}

O orquestrador \texttt{run\_streaming\_pipeline.py} distingue-se dos restantes pela implementação de:

\begin{itemize}[leftmargin=*,noitemsep]
  \item \textbf{Retry com backoff exponencial} --- 3 tentativas automáticas (2s $\rightarrow$ 4s) sobre cada chamada Flyte; erros transitórios de rede não abortam o pipeline
  \item \textbf{Compactação automática Iceberg} --- \texttt{ALTER TABLE ... EXECUTE optimize} após cada run bem-sucedido; elimina o \textit{small-files problem} de ingestões incrementais diárias
  \item \textbf{SLA operacional monitorizado} --- \textit{freshness} máx.\ 7 dias e \textit{runtime} máx.\ 45 min; avisos persistidos em tabela de auditoria
  \item \textbf{Flags de controlo}: \texttt{--source \{energycharts|entsoe\}}, \texttt{--start}, \texttt{--end}, \texttt{--today}, \texttt{--full}, \texttt{--clean}
\end{itemize}

\subsection{Auditoria e Lineage --- DP-02 Streaming}

Cada execução do pipeline Streaming persiste automaticamente dois registos no schema \texttt{iceberg.audit}:

\begin{description}[leftmargin=*]
  \item[\texttt{audit.pipeline\_runs}] Uma linha por run: \texttt{run\_id} (UUID), \texttt{status} (SUCCESS/FAILED), \texttt{duration\_seconds}, contagens de linhas por camada (\texttt{rows\_bronze}, \texttt{rows\_silver}, \texttt{rows\_gold}), parâmetros e mensagem de erro.
  \item[\texttt{audit.dataset\_lineage}] Grafo upstream $\rightarrow$ downstream materializado por execução: \texttt{bronze.consumo\_api\_raw $\rightarrow$ silver.consumo\_api\_hourly $\rightarrow$ gold.dp\_energy\_market\_api\_hourly}.
\end{description}

O logging estruturado prefixo os primeiros 8 caracteres do UUID em cada linha de log, correlacionando diretamente os registos do terminal com as linhas da tabela de auditoria:

\begin{lstlisting}
[3f2a1b9c] Pipeline iniciada — periodo=2024-01-01->2024-12-31  fonte=energycharts
[3f2a1b9c] Bronze ingerido: 17544 linhas  (consumo + preco)
[3f2a1b9c] Pipeline SUCCESS em 541s  (bronze=17544  silver=8772  gold=8748)
\end{lstlisting}

\subsection{Orquestração Flyte}

Os scripts ML são decorados com \texttt{@task} e \texttt{@workflow} do SDK Flyte, com \textit{stubs} locais que permitem execução sem o Flyte instalado:

\begin{lstlisting}[language=Python]
try:
    from flytekit import task, workflow
except ModuleNotFoundError:
    def task(*args, **kwargs):
        def decorator(func): return func
        return args[0] if args and callable(args[0]) else decorator
    def workflow(func=None, **kwargs):
        return func if func else lambda f: f
\end{lstlisting}

O Flyte Sandbox corre como contentor externo com K3s interno, acedendo aos serviços do Compose via \texttt{host.docker.internal}. O DP-02 Static usa execução \textbf{remota} em K3s pods; o DP-02 Streaming usa execução \textbf{local} via \texttt{pyflyte run}.

\subsection{Decisão de Impacto: Reprocessamento Incremental com DELETE+INSERT}

O DP-01, DP-02 Static e DP-03 suportam reprocessamento incremental da camada Gold via \texttt{--date-from} / \texttt{--date-to}. A estratégia DELETE+INSERT (em vez de UPSERT) foi escolhida porque:
\begin{itemize}[leftmargin=*,noitemsep]
  \item Iceberg garante atomicidade --- o DELETE e o INSERT são commitados numa transação
  \item Elimina a necessidade de chave de \textit{merge} e lógica de deduplicação no Gold
  \item Permite re-correr a pipeline para um intervalo sem afetar dados fora desse intervalo
\end{itemize}

% =============================================================================
\section{D) Serving (10\%)}
% =============================================================================

\subsection{Dashboards Grafana como Camada de Serving}

A visualização de todos os Data Products é feita exclusivamente via \textbf{Grafana 10.4}, auto-provisionado a partir de ficheiros JSON em \texttt{04\_application/grafana/dashboards/}. O datasource executa SQL direto sobre Trino/Iceberg --- a Gold é a única fonte de verdade, sem ETL adicional.

\begin{table}[h!]
\centering
\begin{tabularx}{\linewidth}{llX}
\toprule
\textbf{DP} & \textbf{Dashboard} & \textbf{Gold Table consultada} \\
\midrule
DP-01 & \texttt{producao\_consumo\_overview.json} & \texttt{iceberg.gold.dp\_energia\_balance\_hourly} \\
DP-02 Static & \texttt{consumo\_preco\_overview.json} & \texttt{iceberg.gold.dp\_energy\_market\_hourly} \\
DP-02 Streaming & \texttt{consumo\_preco\_streaming\_overview.json} & \texttt{iceberg.gold.dp\_energy\_market\_api\_hourly} \\
DP-03 & \texttt{meteo\_producao\_overview.json} & \texttt{iceberg.gold.dp\_meteo\_producao\_daily\_features} \\
\bottomrule
\end{tabularx}
\end{table}

\subsection{Conteúdo dos Dashboards}

Cada dashboard operacional inclui KPIs em tempo real, séries temporais e distribuições mensais. O DP-02 tem dois dashboards complementares:

\begin{itemize}[leftmargin=*,noitemsep]
  \item \textbf{DP-01}: saldo líquido nacional (MWh), ratio produção/consumo, horas de défice e excedente, série temporal horária
  \item \textbf{DP-02 Static}: preço médio \textit{day-ahead} (€/MWh), custo estimado mensal, correlação preço--consumo, percentagem de horas com preço negativo
  \item \textbf{DP-02 Streaming}: dados quase-em-tempo-real via ENTSO-E, \textit{rolling average} 24h, comparação PT vs.\ ES, indicador de \textit{freshness} ligado à tabela \texttt{audit.pipeline\_runs}
  \item \textbf{DP-03}: temperatura média diária, produção renovável total, \textit{scatter} temperatura vs.\ produção
\end{itemize}

\subsection{Decisão de Impacto: Grafana sobre Trino como Única Camada de Serving}

A decisão de usar Grafana diretamente sobre Trino (sem backend HTTP intermédio) simplifica radicalmente a arquitetura de serving:
\begin{itemize}[leftmargin=*,noitemsep]
  \item Eliminação de três servidores Python (sem risco de \textit{serving skew} entre API e Gold)
  \item Dashboards recarregados via \texttt{POST /api/admin/provisioning/dashboards/reload} sem reiniciar o stack
  \item Provisionamento como código (JSON em git) --- dashboards versionados, reprodutíveis e auditáveis
  \item A Gold Iceberg é a única fonte de verdade para dashboards, ML e qualidade --- sem camadas de cache intermédias
\end{itemize}

% =============================================================================
\section{E) Qualidade e Observabilidade (15\%)}
\label{sec:quality}
% =============================================================================

\subsection{Quality Gates Bloqueantes}

Cada runner executa um \textit{quality gate} automático após a carga da Gold. O padrão de output é uniforme:

\begin{lstlisting}
check_name | status | valor_pct | threshold_pct | detalhe
\end{lstlisting}

\begin{itemize}[leftmargin=*,noitemsep]
  \item \textbf{\textcolor{pass-green}{PASS}}: check dentro dos limites
  \item \textbf{\textcolor{warn-orange}{WARN}}: anomalia não crítica --- pipeline continua, log emitido
  \item \textbf{\textcolor{fail-red}{FAIL}}: violação crítica --- \texttt{SystemExit} imediato; dados não chegam a consumidores
\end{itemize}

Todos os resultados são persistidos automaticamente na tabela \texttt{iceberg.gold.quality\_log} (particionada por \texttt{run\_date}), permitindo auditoria histórica, \textit{trending} de WARNs ao longo do tempo e rastreabilidade de execuções.

\subsection{Cobertura de Checks}

\begin{table}[h!]
\centering
\begin{tabular}{lrrrr}
\toprule
\textbf{DP} & \textbf{Bronze checks} & \textbf{Silver checks} & \textbf{Gold checks} & \textbf{Total} \\
\midrule
DP-01 & --- & 11 & 10 & 21 \\
DP-02 Static & 11 & 13 & 14 & 38 \\
DP-02 Streaming & 10 & 8 & 8 & 26 \\
DP-03 & --- & 13 & 12 & 25 \\
\midrule
\textbf{Total} & & & & \textbf{110} \\
\bottomrule
\end{tabular}
\end{table}

O DP-02 Streaming inclui ainda verificações específicas de \textit{streaming}: \textit{freshness} com limiar de 3 dias (consumo) e 2 dias (preço \textit{day-ahead}), unicidade de \texttt{ts\_utc} na Bronze, e completude diária $\geq 23$ horas (tolerância DST).

Tipos de verificação implementados:

\begin{itemize}[leftmargin=*,noitemsep]
  \item \textbf{Null rates}: colunas críticas devem ter 0\% de nulos (FAIL) ou $<5$\% (WARN)
  \item \textbf{Range físico}: temperatura $\in [-10, 50]$°C, precipitação $\in [0, 200]$~mm, vento $\in [0, 80]$~m/s
  \item \textbf{Unicidade}: chaves de negócio sem duplicados
  \item \textbf{Alinhamento temporal}: \textit{timestamps} alinhados à hora UTC ou a intervalos de 15~min
  \item \textbf{Lag consistency}: \texttt{temp\_lag\_1d} vs temperatura do dia anterior (tolerância 0,01°C)
  \item \textbf{Cobertura cross-DP}: $\geq 90$\% dos dias com meteo + produção + preço
  \item \textbf{Freshness}: dados com atraso $< 400$ dias (threshold adaptado ao dataset histórico)
\end{itemize}

\subsection{Anomalias Esperadas e Conhecimento de Domínio (DP-02)}

Na execução de validação do DP-02 (\textit{consumo\_preco}), todos os 38 checks foram executados: 32 PASS, 6 WARN, 0 FAIL. Os WARNs não bloqueiam o pipeline porque representam características conhecidas do mercado ou artefactos controlados do processo de ingestão:

\begin{table}[h!]
\centering
\small
\begin{tabular}{p{4.8cm}lrp{5.5cm}}
\toprule
\textbf{Check} & \textbf{Camada} & \textbf{Contagem} & \textbf{Explicação} \\
\midrule
\texttt{price\_portugal\_raw >= 0} & Bronze + Silver & 504 horas & Preços negativos são fenómeno real do MIBEL: excesso de produção renovável em períodos de baixa procura empurra o preço abaixo de zero. O dado é válido e preservado; o WARN garante rastreabilidade. \\
\addlinespace
\texttt{uniqueness (datahora,} \texttt{process\_date)} & Bronze & 24 grupos & Artefacto de re-ingest: a Bronze é carregada com idempotência diária, pelo que uma segunda execução sobre o mesmo dia gera duplicados temporários. A Silver absorve-os via \texttt{GROUP BY + AVG}, sem propagação. \\
\addlinespace
\texttt{registos\_por\_dia >= 80} & Bronze & 2 dias & Dias de fronteira do ficheiro histórico com carga parcial ($<96$ registos de 15~min esperados). Não afecta a agregação horária da Silver. \\
\addlinespace
\texttt{horas\_por\_dia >= 23} & Silver & 2 dias & Dias de transição DST (\textit{spring forward}): a hora 2 local não existe, pelo que o dia tem 23 horas. Comportamento esperado e correcto. \\
\bottomrule
\end{tabular}
\caption{WARNs observados no DP-02 e respectiva justificação de domínio.}
\label{tab:warns_dp02}
\end{table}

Esta abordagem --- distinguir anomalias esperadas (WARN) de violações críticas (FAIL) --- é uma decisão de engenharia deliberada: evita \textit{alert fatigue} (falsos alarmes que levam à ignorância de alertas reais) e documenta explicitamente o comportamento do mercado energético ibérico no próprio sistema de qualidade.

\subsection{Dashboards Grafana de Observabilidade}

Quatro dashboards operacionais provisionados automaticamente (um por DP/variante) --- ver Secção~D. Adicionalmente, o DP-02 Streaming expõe métricas de auditoria diretamente no dashboard via queries sobre \texttt{iceberg.audit.pipeline\_runs}:

\begin{itemize}[leftmargin=*,noitemsep]
  \item Histograma de \textit{duration\_seconds} por execução --- deteta degradação de performance
  \item Série temporal de \texttt{rows\_gold} por run --- deteta gaps de ingestão
  \item Painel de \textit{freshness}: \texttt{CURRENT\_DATE - MAX(process\_date)} em dias
  \item Tabela das últimas 10 execuções com status, duração e contagens de linhas
\end{itemize}

Esta integração fecha o ciclo de observabilidade: o mesmo pipeline que persiste os dados de negócio persiste também os metadados de execução, visíveis no mesmo Grafana sem infraestrutura adicional.

\subsection{Decisão de Impacto: Quality Gate Bloqueante vs. Apenas Alertar}

A decisão de bloquear o pipeline em FAIL (em vez de apenas emitir um alerta) tem impacto direto na qualidade downstream:

\begin{itemize}[leftmargin=*,noitemsep]
  \item Um erro de qualidade na Gold propagaria para modelos ML (métricas inflacionadas), dashboards (valores incorretos) e APIs (respostas erradas)
  \item O custo de um FAIL na pipeline é baixo (re-correr após correção)
  \item O custo de um modelo ML treinado com dados corrompidos é alto (re-treino, validação, re-deployment)
\end{itemize}

% =============================================================================
\section{F) Machine Learning (15\%)}
% =============================================================================

\subsection{Modelos e Feature Tables}

\begin{table}[h!]
\centering
\begin{tabularx}{\linewidth}{lllX}
\toprule
\textbf{DP} & \textbf{Modelo} & \textbf{Target} & \textbf{Feature Table (Gold)} \\
\midrule
DP-01 & RandomForestClassifier & \texttt{flag\_defice} (t+1h) & \texttt{dp\_energia\_balance\_hourly} \\
DP-02 Static & GradientBoostingRegressor & \texttt{consumo\_next\_hour} & \texttt{feat\_load\_forecasting\_hourly} \\
DP-02 Streaming & GradientBoostingRegressor & \texttt{consumo\_next\_hour} & \texttt{feat\_load\_forecasting\_api\_hourly} \\
DP-03 A & RandomForestRegressor & \texttt{producao\_total\_daily\_mwh} & \texttt{dp\_meteo\_producao\_daily\_features} \\
DP-03 B & GradientBoostingRegressor & \texttt{preco\_spot\_medio\_eur\_mwh} & \texttt{dp\_meteo\_producao\_daily\_features} \\
\bottomrule
\end{tabularx}
\end{table}

O DP-02 Streaming tem a sua própria feature table (\texttt{feat\_load\_forecasting\_api\_hourly}) derivada dos dados ENTSO-E, com as mesmas colunas que a variante estática mas particionada por \texttt{year/month} para eficiência em ingestões incrementais.

\subsection{Rastreabilidade de Dados por Versão de Modelo}

Cada \textit{run} MLflow regista:

\begin{itemize}[leftmargin=*,noitemsep]
  \item \textbf{Tags}: \texttt{dataset} (tabela Gold exata), \texttt{target}, \texttt{model\_type}, \texttt{domain}
  \item \textbf{Params}: hiperparâmetros + \texttt{train\_start/end}, \texttt{test\_start/end}, \texttt{n\_train}, \texttt{n\_test}, \texttt{random\_state}
  \item \textbf{Métricas}: avaliação no conjunto de teste (MAE, RMSE, R², MAPE / Accuracy, F1, ROC-AUC)
  \item \textbf{Artefactos}: \texttt{feature\_importances.json}, \texttt{classification\_report.json}, \texttt{feature\_columns.json}, \texttt{model/} (pipeline sklearn serializado)
  \item \textbf{Model Registry}: 4 modelos com \texttt{registered\_model\_name}
\end{itemize}

\subsection{Decisão de Impacto: Split Temporal e Feature Tables Pré-computadas}

\textbf{Split temporal sem shuffle}: séries temporais têm dependência sequencial; embaralhar os dados introduz \textit{data leakage} (o modelo veria o futuro durante o treino), inflacionando as métricas de avaliação. O split temporal garante que o conjunto de teste representa dados genuinamente não vistos.

\textbf{Feature tables pré-computadas na Gold}: lags e \textit{rolling windows} são calculados uma única vez em SQL Trino e armazenados em Iceberg. As alternativas (calcular em Python durante o treino) introduzem risco de \textit{training-serving skew} --- se o cálculo em produção diferir minimamente do cálculo em treino, as previsões degradam silenciosamente.

% =============================================================================
\section{Relevância para o Negócio}
% =============================================================================

\subsection{Impacto Hipotético dos Três Data Products}

\subsubsection{DP-01 — Monitorização de Défice Energético}

Um comercializador que preveja défice com 1 hora de antecedência pode antecipar compras no mercado intradiário a preços mais favoráveis do que no mercado de balanço (tipicamente mais caro e menos previsível). O modelo de classificação (RF, target \texttt{flag\_defice} t+1h) oferece directamente este sinal de decisão.

\subsubsection{DP-02 — Estimativa de Custo Horário}

A feature table \texttt{feat\_load\_forecasting\_hourly} com lags 1h/24h e \textit{rolling} 24h permite ao modelo GB prever o consumo da próxima hora. Num cenário real, esta previsão alimentaria um sistema de otimização de portefólio: saber se o consumo sobe ou desce na próxima hora determina se é vantajoso comprar antecipadamente ou aguardar.

\subsubsection{DP-03 — Previsão Meteo $\rightarrow$ Produção $\rightarrow$ Preço}

A cadeia de modelos (Modelo A: produção a partir de meteorologia; Modelo B: preço a partir de produção + meteorologia) reflecte uma hipótese de negócio real: em mercados com elevada penetração renovável, a produção fotovoltaica/eólica (altamente dependente do tempo) é o principal determinante do preço \textit{spot}. Um modelo que antecipe o preço a partir da previsão meteorológica (disponível 7 dias antes) permite estratégias de \textit{hedging} com muito maior antecedência.

\subsection{Valor da Qualidade como Diferenciador}

Os 84 checks de qualidade e os dashboards de observabilidade não são apenas boas práticas técnicas --- têm impacto direto em decisões de negócio. Um sinal de défice gerado a partir de dados com duplicados ou valores fora de range pode levar a uma decisão de compra desnecessária, com custo financeiro real. A rastreabilidade dados $\leftrightarrow$ modelo (via MLflow tags e params) permite auditar \textit{a posteriori} qualquer decisão automatizada que tenha usado um modelo específico.

% =============================================================================
\section{Conclusão}
% =============================================================================

O projeto MIBEL implementou integralmente os seis critérios de avaliação:

\begin{itemize}[leftmargin=*]
  \item \textbf{A)} Três Data Products especificados com contratos formais, SLAs e estratégia de versionamento
  \item \textbf{B)} Lakehouse local completo com Bronze/Silver/Gold em Iceberg + Trino sobre MinIO
  \item \textbf{C)} Pipelines idempotentes com quality gates, reprocessamento incremental, orquestração Flyte e --- no DP-02 Streaming --- auditoria completa, retry com backoff exponencial e compactação automática Iceberg
  \item \textbf{D)} Quatro dashboards Grafana provisionados como código (JSON em git), com SQL direto sobre Trino --- sem camadas de serving intermédias
  \item \textbf{E)} 110 checks de qualidade distribuídos por 4 pipelines (DP-01, DP-02 Static, DP-02 Streaming, DP-03); tabelas de auditoria \texttt{iceberg.audit} com lineage materializado por execução
  \item \textbf{F)} 5 modelos ML com rastreabilidade completa de dados, split temporal e feature tables pré-computadas na Gold
\end{itemize}

As decisões de maior impacto foram: separação Hive/Iceberg por camada, quality gate bloqueante em FAIL, Grafana sobre Trino como única camada de serving, split temporal nos modelos ML, feature tables pré-computadas na Gold e auditoria \texttt{iceberg.audit} com lineage automático no DP-02 Streaming. Todas estas decisões garantem que dados, modelos e execuções em produção reflectem fielmente a realidade do mercado energético português e são auditáveis \textit{a posteriori}.

% =============================================================================
\end{document}
